Material Acceleration

A moving particle can accelerate even if the field is steady.

MATERIAL DERIVATIVE
$$\frac{D}{Dt}=\frac{\partial}{\partial t}$$
$$+u\frac{\partial}{\partial x}+v\frac{\partial}{\partial y}$$
$$+w\frac{\partial}{\partial z}$$

X ACCELERATION
$$a_x=u_t+uu_x+vu_y+wu_z$$
Y ACCELERATION
$$a_y=v_t+uv_x+vv_y+wv_z$$
Z ACCELERATION
$$a_z=w_t+uw_x+vw_y+ww_z$$
Subscript derivative e.g. $u_x=\partial u/\partial x$.

LOCAL = partial-time term.
CONVECTIVE = velocity times spatial gradients.

PROCEDURE
1. Make derivative table: t, x, y, z.
2. Form local and convective parts separately.
3. Add; THEN substitute point coordinates.
4. Magnitude: $|a|=\sqrt{a_x^2+a_y^2+a_z^2}$.
5. Check quadrant of direction (atan2).

EXAMPLE
$u=-Ax+Bt$; $v=Ay+Bt$; w=0.
Local accel = (B,B,0).
Convective x = $A^2x-ABt$.
Convective y = $A^2y+ABt$.
For A=1, B=2, x=1, y=2, t=5:
a = (-7,14) m/s^2.
Magnitude 15.65; angle 116.57 deg from +x.
